% ============================================
% Section: Conclusion
% ============================================
\label{sec:conclusion}

We presented Co-LIC2, a cooperative SLAM system that extends single-platform
LiDAR-Inertial-Visual Gaussian Splatting SLAM to a team of communicating
moving platforms. The system generalizes the continuous-time multi-sensor
fusion front-end of Coco-LIC and Gaussian-LIC2 to an arbitrary number of
cameras and LiDARs per node, introduces a hierarchical hybrid communication
architecture with cluster-head election and a bandwidth-aware submap
publication protocol, and reconciles per-node Gaussian maps into a globally
consistent photo-realistic map through Gaussian-to-Gaussian ICP registration,
ownership-tagged confidence-weighted merging, and robust pose-graph
optimization. A failure-mode analysis shows the system degrades gracefully to
single-node Gaussian-LIC2 under total partition, and a bandwidth model
demonstrates sub-linear per-link load growth with node count.

The design is fully specified in a companion design document, and the
implementation reuses the open-source Gaussian-LIC2 code base, preserving
single-node behavior when the cooperative layer is disabled. An empirical
evaluation on physical multi-rig hardware is the subject of ongoing work; this
report fixes the system model and algorithms against which those results will
be read. We believe Co-LIC2 is a concrete step toward practical cooperative
photo-realistic mapping with heterogeneous multi-sensor teams.
